%%%%%%%%%%%%%%%%%%%%%%%%%%%%%%%%%%%%%%%%%%%%%%%%%%%%%%%%%%%%
% 7.15 RANDOM MATRIX THEORY
% The Composition of Advanced Mathematics
%%%%%%%%%%%%%%%%%%%%%%%%%%%%%%%%%%%%%%%%%%%%%%%%%%%%%%%%%%%%

\section{Random Matrix Theory}
\label{sec:random-matrix-theory}

\prerequisite{
Probability distributions, random variables, expectation, variance,
joint distributions, multivariate probability, linear algebra,
eigenvalues, eigenvectors, symmetric and Hermitian matrices,
determinants, traces, matrix norms, asymptotic probability, and basic
measure-theoretic probability.
}

\learningobjectives{
By the end of this section, the reader should be able to:
\begin{itemize}
    \item define a random matrix;
    \item distinguish entrywise randomness from spectral randomness;
    \item define empirical spectral distributions;
    \item compute trace and determinant statistics;
    \item understand Wigner matrices;
    \item state the semicircle law;
    \item understand normalized eigenvalue scaling;
    \item define Gaussian orthogonal and unitary ensembles;
    \item interpret eigenvalue repulsion;
    \item understand Vandermonde factors in eigenvalue densities;
    \item define sample covariance matrices;
    \item define Wishart matrices;
    \item state the Marchenko--Pastur law;
    \item understand aspect ratios in high-dimensional covariance;
    \item recognize the largest-eigenvalue problem;
    \item understand the Tracy--Widom distribution conceptually;
    \item distinguish bulk and edge spectral behavior;
    \item understand singular values of random matrices;
    \item connect random matrices with principal component analysis;
    \item understand spiked covariance models;
    \item recognize the BBP phase transition conceptually;
    \item understand random graph adjacency matrices as random matrices;
    \item interpret random matrix universality;
    \item understand concentration of spectral statistics;
    \item recognize resolvents and Stieltjes transforms;
    \item connect random matrix theory with statistics, physics,
          numerical mathematics, networks, machine learning, and number
          theory.
\end{itemize}
}

%%%%%%%%%%%%%%%%%%%%%%%%%%%%%%%%%%%%%%%%%%%%%%%%%%%%%%%%%%%%
\subsection{What Is a Random Matrix?}
\label{subsec:what-is-random-matrix}
%%%%%%%%%%%%%%%%%%%%%%%%%%%%%%%%%%%%%%%%%%%%%%%%%%%%%%%%%%%%

A random matrix is a matrix whose entries are random variables.

\begin{definitionbox}{Random Matrix}
A random matrix is a matrix-valued random variable

\[
\boxed{
A:\Omega\to\mathbb{R}^{m\times n}
}
\]

or, more generally,

\[
\boxed{
A:\Omega\to\mathbb{C}^{m\times n}.
}
\]
\end{definitionbox}

Thus each outcome \(\omega\) produces an ordinary matrix

\[
A(\omega).
\]

The randomness may appear in every entry or only in part of the matrix.

%%%%%%%%%%%%%%%%%%%%%%%%%%%%%%%%%%%%%%%%%%%%%%%%%%%%%%%%%%%%
\subsection{Why Random Matrices Matter}
\label{subsec:why-random-matrices}
%%%%%%%%%%%%%%%%%%%%%%%%%%%%%%%%%%%%%%%%%%%%%%%%%%%%%%%%%%%%

Random matrices arise whenever a large system contains many random or
uncertain interactions.

Examples include:

\[
\boxed{
\text{sample covariance matrices},
}
\]

\[
\boxed{
\text{random networks},
}
\]

\[
\boxed{
\text{quantum systems},
}
\]

\[
\boxed{
\text{high-dimensional data},
}
\]

\[
\boxed{
\text{wireless communication channels},
}
\]

and

\[
\boxed{
\text{random numerical algorithms}.
}
\]

A central goal is to understand the eigenvalues and singular values of
large random matrices.

%%%%%%%%%%%%%%%%%%%%%%%%%%%%%%%%%%%%%%%%%%%%%%%%%%%%%%%%%%%%
\subsection{Entrywise Versus Spectral Questions}
\label{subsec:entrywise-versus-spectral}
%%%%%%%%%%%%%%%%%%%%%%%%%%%%%%%%%%%%%%%%%%%%%%%%%%%%%%%%%%%%

An entrywise question asks about quantities such as

\[
A_{ij}.
\]

A spectral question asks about quantities such as

\[
\lambda_1,\ldots,\lambda_n,
\]

where the \(\lambda_i\) are eigenvalues.

Random matrix theory is primarily spectral.

Thus the main transition is

\[
\boxed{
\text{random entries}
\longrightarrow
\text{random eigenvalues}.
}
\]

%%%%%%%%%%%%%%%%%%%%%%%%%%%%%%%%%%%%%%%%%%%%%%%%%%%%%%%%%%%%
\subsection{Basic Random Matrix Model}
\label{subsec:basic-random-matrix-model}
%%%%%%%%%%%%%%%%%%%%%%%%%%%%%%%%%%%%%%%%%%%%%%%%%%%%%%%%%%%%

Consider an \(n\times n\) matrix

\[
A
=
(A_{ij}),
\]

where the entries are independent random variables.

For example,

\[
A_{ij}
\sim
N(0,1).
\]

This defines a Gaussian random matrix.

However, if \(A\) is not symmetric, its eigenvalues may be complex.

To study real eigenvalues, symmetric models are especially important.

%%%%%%%%%%%%%%%%%%%%%%%%%%%%%%%%%%%%%%%%%%%%%%%%%%%%%%%%%%%%
\subsection{Random Symmetric Matrices}
\label{subsec:random-symmetric-matrices}
%%%%%%%%%%%%%%%%%%%%%%%%%%%%%%%%%%%%%%%%%%%%%%%%%%%%%%%%%%%%

A real symmetric random matrix satisfies

\[
\boxed{
A=A^T.
}
\]

Therefore,

\[
A_{ij}=A_{ji}.
\]

Only the entries on and above the diagonal need to be generated
independently.

All eigenvalues are real.

This makes symmetric random matrices natural objects for spectral
analysis.

%%%%%%%%%%%%%%%%%%%%%%%%%%%%%%%%%%%%%%%%%%%%%%%%%%%%%%%%%%%%
\subsection{Hermitian Random Matrices}
\label{subsec:hermitian-random-matrices}
%%%%%%%%%%%%%%%%%%%%%%%%%%%%%%%%%%%%%%%%%%%%%%%%%%%%%%%%%%%%

In the complex setting, the analogue of symmetry is Hermitian
structure:

\[
\boxed{
A=A^*,
}
\]

where

\[
A^*
=
\overline{A}^{\,T}.
\]

Hermitian matrices also have real eigenvalues.

Many important random matrix ensembles are Hermitian.

%%%%%%%%%%%%%%%%%%%%%%%%%%%%%%%%%%%%%%%%%%%%%%%%%%%%%%%%%%%%
\subsection{Spectral Decomposition}
\label{subsec:random-spectral-decomposition}
%%%%%%%%%%%%%%%%%%%%%%%%%%%%%%%%%%%%%%%%%%%%%%%%%%%%%%%%%%%%

If \(A\) is real symmetric, then

\[
\boxed{
A
=
Q\Lambda Q^T,
}
\]

where \(Q\) is orthogonal and

\[
\Lambda
=
\operatorname{diag}
(
\lambda_1,\ldots,\lambda_n
).
\]

The eigenvalues

\[
\lambda_1,\ldots,\lambda_n
\]

are themselves random variables.

Thus a random matrix produces a random spectrum.

%%%%%%%%%%%%%%%%%%%%%%%%%%%%%%%%%%%%%%%%%%%%%%%%%%%%%%%%%%%%
\subsection{Empirical Spectral Distribution}
\label{subsec:empirical-spectral-distribution}
%%%%%%%%%%%%%%%%%%%%%%%%%%%%%%%%%%%%%%%%%%%%%%%%%%%%%%%%%%%%

Suppose an \(n\times n\) Hermitian matrix has eigenvalues

\[
\lambda_1,\ldots,\lambda_n.
\]

\begin{definitionbox}{Empirical Spectral Distribution}
The empirical spectral distribution, abbreviated ESD, is the probability
measure

\[
\boxed{
\mu_A
=
\frac1n
\sum_{i=1}^{n}
\delta_{\lambda_i},
}
\]

where \(\delta_{\lambda_i}\) denotes a point mass at \(\lambda_i\).
\end{definitionbox}

The ESD places equal mass \(1/n\) at every eigenvalue.

%%%%%%%%%%%%%%%%%%%%%%%%%%%%%%%%%%%%%%%%%%%%%%%%%%%%%%%%%%%%
\subsection{Spectral Distribution Function}
\label{subsec:spectral-distribution-function}
%%%%%%%%%%%%%%%%%%%%%%%%%%%%%%%%%%%%%%%%%%%%%%%%%%%%%%%%%%%%

The corresponding cumulative spectral distribution is

\[
\boxed{
F_A(x)
=
\frac1n
\#\{i:\lambda_i\leq x\}.
}
\]

Thus \(F_A(x)\) is the fraction of eigenvalues at or below \(x\).

For large \(n\), random matrix theory studies whether

\[
F_A(x)
\]

approaches a deterministic limiting distribution.

%%%%%%%%%%%%%%%%%%%%%%%%%%%%%%%%%%%%%%%%%%%%%%%%%%%%%%%%%%%%
\subsection{Trace and Eigenvalues}
\label{subsec:trace-random-matrices}
%%%%%%%%%%%%%%%%%%%%%%%%%%%%%%%%%%%%%%%%%%%%%%%%%%%%%%%%%%%%

For any square matrix,

\[
\boxed{
\operatorname{tr}(A)
=
\sum_{i=1}^{n}\lambda_i.
}
\]

Thus the trace is a linear spectral statistic.

More generally,

\[
\boxed{
\operatorname{tr}(A^k)
=
\sum_{i=1}^{n}\lambda_i^k.
}
\]

This identity is fundamental in the moment method.

%%%%%%%%%%%%%%%%%%%%%%%%%%%%%%%%%%%%%%%%%%%%%%%%%%%%%%%%%%%%
\subsection{Determinant and Eigenvalues}
\label{subsec:determinant-random-matrices}
%%%%%%%%%%%%%%%%%%%%%%%%%%%%%%%%%%%%%%%%%%%%%%%%%%%%%%%%%%%%

The determinant satisfies

\[
\boxed{
\det(A)
=
\prod_{i=1}^{n}\lambda_i.
}
\]

Therefore,

\[
\log|\det(A)|
=
\sum_{i=1}^{n}
\log|\lambda_i|
\]

whenever the eigenvalues are nonzero.

Determinants therefore encode multiplicative spectral behavior.

%%%%%%%%%%%%%%%%%%%%%%%%%%%%%%%%%%%%%%%%%%%%%%%%%%%%%%%%%%%%
\subsection{Spectral Norm}
\label{subsec:spectral-norm-random}
%%%%%%%%%%%%%%%%%%%%%%%%%%%%%%%%%%%%%%%%%%%%%%%%%%%%%%%%%%%%

For a symmetric matrix,

\[
\boxed{
\|A\|_2
=
\max_i|\lambda_i|.
}
\]

Thus the matrix spectral norm is governed by the extreme eigenvalues.

For rectangular matrices,

\[
\|A\|_2
\]

is the largest singular value.

%%%%%%%%%%%%%%%%%%%%%%%%%%%%%%%%%%%%%%%%%%%%%%%%%%%%%%%%%%%%
\subsection{Wigner Matrices}
\label{subsec:wigner-matrices}
%%%%%%%%%%%%%%%%%%%%%%%%%%%%%%%%%%%%%%%%%%%%%%%%%%%%%%%%%%%%

Wigner matrices are foundational models in random matrix theory.

\begin{definitionbox}{Wigner Matrix}
A real Wigner matrix is a symmetric matrix

\[
W_n=(W_{ij})_{1\leq i,j\leq n}
\]

whose upper-triangular entries are independent, typically with

\[
\mathbb{E}[W_{ij}]=0
\]

and

\[
\operatorname{Var}(W_{ij})=\sigma^2
\]

for \(i<j\), together with suitable moment assumptions.
\end{definitionbox}

The symmetry condition imposes

\[
W_{ji}=W_{ij}.
\]

%%%%%%%%%%%%%%%%%%%%%%%%%%%%%%%%%%%%%%%%%%%%%%%%%%%%%%%%%%%%
\subsection{Why Wigner Matrices Must Be Rescaled}
\label{subsec:wigner-scaling}
%%%%%%%%%%%%%%%%%%%%%%%%%%%%%%%%%%%%%%%%%%%%%%%%%%%%%%%%%%%%

If the entries have variance of order \(1\), then the eigenvalues of
\(W_n\) are typically of order

\[
\sqrt{n}.
\]

Therefore one studies the normalized matrix

\[
\boxed{
M_n
=
\frac{1}{\sqrt{n}}W_n.
}
\]

This scaling keeps the bulk spectrum on an order-one interval as

\[
n\to\infty.
\]

%%%%%%%%%%%%%%%%%%%%%%%%%%%%%%%%%%%%%%%%%%%%%%%%%%%%%%%%%%%%
\subsection{Wigner Semicircle Law}
\label{subsec:wigner-semicircle-law}
%%%%%%%%%%%%%%%%%%%%%%%%%%%%%%%%%%%%%%%%%%%%%%%%%%%%%%%%%%%%

\begin{theorem}[Wigner Semicircle Law]
Under standard Wigner assumptions with off-diagonal variance
\(\sigma^2\), the empirical spectral distribution of

\[
\frac1{\sqrt{n}}W_n
\]

converges, in the appropriate probabilistic sense, to the semicircle
distribution with density

\[
\boxed{
\rho_{\mathrm{sc}}(x)
=
\frac{
1
}{
2\pi\sigma^2
}
\sqrt{
4\sigma^2-x^2
}
}
\]

for

\[
|x|\leq2\sigma,
\]

and

\[
\rho_{\mathrm{sc}}(x)=0
\]

otherwise.
\end{theorem}

%%%%%%%%%%%%%%%%%%%%%%%%%%%%%%%%%%%%%%%%%%%%%%%%%%%%%%%%%%%%
\subsection{Standard Semicircle Density}
\label{subsec:standard-semicircle-density}
%%%%%%%%%%%%%%%%%%%%%%%%%%%%%%%%%%%%%%%%%%%%%%%%%%%%%%%%%%%%

When

\[
\sigma^2=1,
\]

the density becomes

\[
\boxed{
\rho_{\mathrm{sc}}(x)
=
\frac1{2\pi}
\sqrt{
4-x^2
},
\qquad
-2\leq x\leq2.
}
\]

The graph forms the upper half of a semicircular profile after
appropriate scaling.

%%%%%%%%%%%%%%%%%%%%%%%%%%%%%%%%%%%%%%%%%%%%%%%%%%%%%%%%%%%%
\subsection{Interpretation of the Semicircle Law}
\label{subsec:interpret-semicircle-law}
%%%%%%%%%%%%%%%%%%%%%%%%%%%%%%%%%%%%%%%%%%%%%%%%%%%%%%%%%%%%

The semicircle law does not say that every eigenvalue is deterministic.

Instead, it says that the overall histogram of many eigenvalues becomes
predictable.

Thus:

\[
\boxed{
\text{individual eigenvalues remain random}
}
\]

while

\[
\boxed{
\text{global spectral density becomes deterministic}.
}
\]

This is a recurring theme in large random systems.

%%%%%%%%%%%%%%%%%%%%%%%%%%%%%%%%%%%%%%%%%%%%%%%%%%%%%%%%%%%%
\subsection{Worked Example: Trace Scaling}
\label{subsec:worked-trace-wigner}
%%%%%%%%%%%%%%%%%%%%%%%%%%%%%%%%%%%%%%%%%%%%%%%%%%%%%%%%%%%%

\begin{workedexamplebox}{Normalized Trace of a Centered Wigner Matrix}

Let

\[
M_n
=
\frac1{\sqrt{n}}W_n
\]

where the diagonal entries have mean zero.

Then

\[
\operatorname{tr}(M_n)
=
\frac1{\sqrt{n}}
\sum_{i=1}^{n}W_{ii}.
\]

The normalized first empirical spectral moment is

\[
\frac1n
\operatorname{tr}(M_n).
\]

Under standard moment assumptions this tends toward zero.

\begin{answerbox}
\[
\boxed{
\int x\,d\mu_{M_n}(x)
=
\frac1n\operatorname{tr}(M_n)
\to0,
}
\]

consistent with the symmetry of the semicircle law.
\end{answerbox}

\end{workedexamplebox}

%%%%%%%%%%%%%%%%%%%%%%%%%%%%%%%%%%%%%%%%%%%%%%%%%%%%%%%%%%%%
\subsection{Moment Method}
\label{subsec:random-matrix-moment-method}
%%%%%%%%%%%%%%%%%%%%%%%%%%%%%%%%%%%%%%%%%%%%%%%%%%%%%%%%%%%%

One way to prove spectral convergence is to study moments.

For the ESD,

\[
\int x^k\,d\mu_A(x)
=
\frac1n
\sum_{i=1}^{n}\lambda_i^k.
\]

Using the trace identity,

\[
\boxed{
\int x^k\,d\mu_A(x)
=
\frac1n
\operatorname{tr}(A^k).
}
\]

Thus spectral moments can be studied through matrix products.

%%%%%%%%%%%%%%%%%%%%%%%%%%%%%%%%%%%%%%%%%%%%%%%%%%%%%%%%%%%%
\subsection{Expanding the Trace Moment}
\label{subsec:trace-moment-expansion}
%%%%%%%%%%%%%%%%%%%%%%%%%%%%%%%%%%%%%%%%%%%%%%%%%%%%%%%%%%%%

For an \(n\times n\) matrix,

\[
\operatorname{tr}(A^k)
=
\sum_{i_1,\ldots,i_k}
A_{i_1i_2}
A_{i_2i_3}
\cdots
A_{i_ki_1}.
\]

Taking expectation gives

\[
\boxed{
\mathbb{E}
[
\operatorname{tr}(A^k)
]
=
\sum_{i_1,\ldots,i_k}
\mathbb{E}
[
A_{i_1i_2}\cdots A_{i_ki_1}
].
}
\]

Independence and centering force many terms to vanish.

%%%%%%%%%%%%%%%%%%%%%%%%%%%%%%%%%%%%%%%%%%%%%%%%%%%%%%%%%%%%
\subsection{Catalan Numbers and the Semicircle Law}
\label{subsec:catalan-semicircle}
%%%%%%%%%%%%%%%%%%%%%%%%%%%%%%%%%%%%%%%%%%%%%%%%%%%%%%%%%%%%

The even moments of the standard semicircle distribution are Catalan
numbers.

For

\[
k\geq0,
\]

\[
\boxed{
\int_{-2}^{2}
x^{2k}
\rho_{\mathrm{sc}}(x)\,dx
=
C_k,
}
\]

where

\[
\boxed{
C_k
=
\frac1{k+1}
\binom{2k}{k}.
}
\]

Odd moments vanish.

This produces an important connection between random matrices and
enumerative combinatorics.

%%%%%%%%%%%%%%%%%%%%%%%%%%%%%%%%%%%%%%%%%%%%%%%%%%%%%%%%%%%%
\subsection{Gaussian Orthogonal Ensemble}
\label{subsec:goe}
%%%%%%%%%%%%%%%%%%%%%%%%%%%%%%%%%%%%%%%%%%%%%%%%%%%%%%%%%%%%

The Gaussian Orthogonal Ensemble, abbreviated GOE, is a probability
distribution on real symmetric matrices.

A common normalization chooses entries so that the matrix density is
proportional to

\[
\boxed{
\exp
\left(
-\frac{n}{4}
\operatorname{tr}(A^2)
\right).
}
\]

The distribution is invariant under orthogonal conjugation:

\[
\boxed{
A
\overset{d}{=}
Q A Q^T
}
\]

for fixed orthogonal \(Q\).

%%%%%%%%%%%%%%%%%%%%%%%%%%%%%%%%%%%%%%%%%%%%%%%%%%%%%%%%%%%%
\subsection{Gaussian Unitary Ensemble}
\label{subsec:gue}
%%%%%%%%%%%%%%%%%%%%%%%%%%%%%%%%%%%%%%%%%%%%%%%%%%%%%%%%%%%%

The Gaussian Unitary Ensemble, abbreviated GUE, consists of Hermitian
random matrices with a distribution invariant under unitary
conjugation.

Its matrix density is of the general form

\[
\boxed{
\exp
\left(
-\frac{n}{2}
\operatorname{tr}(A^2)
\right)
}
\]

up to normalization conventions.

The exact constants depend on the chosen scaling.

%%%%%%%%%%%%%%%%%%%%%%%%%%%%%%%%%%%%%%%%%%%%%%%%%%%%%%%%%%%%
\subsection{Gaussian Symplectic Ensemble}
\label{subsec:gse}
%%%%%%%%%%%%%%%%%%%%%%%%%%%%%%%%%%%%%%%%%%%%%%%%%%%%%%%%%%%%

A third classical Gaussian ensemble is the Gaussian Symplectic
Ensemble, abbreviated GSE.

Together, GOE, GUE, and GSE correspond to three major symmetry classes.

They are often described through the parameter

\[
\boxed{
\beta
=
1,2,4,
}
\]

respectively.

%%%%%%%%%%%%%%%%%%%%%%%%%%%%%%%%%%%%%%%%%%%%%%%%%%%%%%%%%%%%
\subsection{Beta Ensembles}
\label{subsec:beta-ensembles}
%%%%%%%%%%%%%%%%%%%%%%%%%%%%%%%%%%%%%%%%%%%%%%%%%%%%%%%%%%%%

The parameter \(\beta\) can be generalized beyond

\[
1,2,4.
\]

A beta ensemble has eigenvalue density of the form

\[
\boxed{
p(\lambda_1,\ldots,\lambda_n)
\propto
\prod_{i<j}
|\lambda_i-\lambda_j|^\beta
\prod_{i=1}^{n}
e^{-nV(\lambda_i)},
}
\]

for an appropriate potential \(V\).

The factor

\[
\prod_{i<j}
|\lambda_i-\lambda_j|^\beta
\]

controls eigenvalue interaction.

%%%%%%%%%%%%%%%%%%%%%%%%%%%%%%%%%%%%%%%%%%%%%%%%%%%%%%%%%%%%
\subsection{Vandermonde Determinant}
\label{subsec:vandermonde-random-matrix}
%%%%%%%%%%%%%%%%%%%%%%%%%%%%%%%%%%%%%%%%%%%%%%%%%%%%%%%%%%%%

The Vandermonde determinant is

\[
\boxed{
\Delta(\lambda_1,\ldots,\lambda_n)
=
\prod_{i<j}
(\lambda_j-\lambda_i).
}
\]

Random matrix eigenvalue densities often contain

\[
\boxed{
|\Delta(\lambda)|^\beta.
}
\]

This factor strongly suppresses configurations where eigenvalues are
very close.

%%%%%%%%%%%%%%%%%%%%%%%%%%%%%%%%%%%%%%%%%%%%%%%%%%%%%%%%%%%%
\subsection{Eigenvalue Repulsion}
\label{subsec:eigenvalue-repulsion}
%%%%%%%%%%%%%%%%%%%%%%%%%%%%%%%%%%%%%%%%%%%%%%%%%%%%%%%%%%%%

Because of the factor

\[
|\lambda_i-\lambda_j|^\beta,
\]

the probability density tends to zero when

\[
\lambda_i\to\lambda_j.
\]

Thus eigenvalues tend to repel one another.

This phenomenon is called eigenvalue repulsion.

It is fundamentally different from a collection of independent random
points.

%%%%%%%%%%%%%%%%%%%%%%%%%%%%%%%%%%%%%%%%%%%%%%%%%%%%%%%%%%%%
\subsection{Bulk and Edge}
\label{subsec:bulk-and-edge-rmt}
%%%%%%%%%%%%%%%%%%%%%%%%%%%%%%%%%%%%%%%%%%%%%%%%%%%%%%%%%%%%

Random matrix spectra are often studied in two regimes.

The \textbf{bulk} refers to eigenvalues in the interior of the limiting
spectral support.

The \textbf{edge} refers to eigenvalues near the largest or smallest
spectral values.

Thus:

\[
\boxed{
\text{bulk}
=
\text{interior spectral behavior},
}
\]

\[
\boxed{
\text{edge}
=
\text{extreme spectral behavior}.
}
\]

Different limiting laws appear in these regions.

%%%%%%%%%%%%%%%%%%%%%%%%%%%%%%%%%%%%%%%%%%%%%%%%%%%%%%%%%%%%
\subsection{Largest Eigenvalue}
\label{subsec:largest-eigenvalue-rmt}
%%%%%%%%%%%%%%%%%%%%%%%%%%%%%%%%%%%%%%%%%%%%%%%%%%%%%%%%%%%%

For a normalized Wigner matrix, the largest eigenvalue typically lies
near the upper edge of the semicircle support.

Under standard normalization,

\[
\boxed{
\lambda_{\max}
\approx
2
}
\]

for large \(n\).

The deviations from this edge occur on a much smaller scale than the
overall spectral width.

%%%%%%%%%%%%%%%%%%%%%%%%%%%%%%%%%%%%%%%%%%%%%%%%%%%%%%%%%%%%
\subsection{Tracy--Widom Distribution Preview}
\label{subsec:tracy-widom-preview}
%%%%%%%%%%%%%%%%%%%%%%%%%%%%%%%%%%%%%%%%%%%%%%%%%%%%%%%%%%%%

For several classical random matrix ensembles, the centered and scaled
largest eigenvalue converges to a Tracy--Widom distribution.

Symbolically,

\[
\boxed{
n^{2/3}
(
\lambda_{\max}-\lambda_{\mathrm{edge}}
)
\xrightarrow{d}
\text{Tracy--Widom}.
}
\]

The exact scaling constants depend on the ensemble and normalization.

The Tracy--Widom laws describe universal edge fluctuations.

%%%%%%%%%%%%%%%%%%%%%%%%%%%%%%%%%%%%%%%%%%%%%%%%%%%%%%%%%%%%
\subsection{Bulk Spacing Preview}
\label{subsec:bulk-spacing-preview}
%%%%%%%%%%%%%%%%%%%%%%%%%%%%%%%%%%%%%%%%%%%%%%%%%%%%%%%%%%%%

After appropriate local rescaling, neighboring eigenvalue spacings in
the bulk exhibit universal behavior.

For the GUE, local correlations are described by the sine kernel:

\[
\boxed{
K(x,y)
=
\frac{
\sin\pi(x-y)
}{
\pi(x-y)
}.
}
\]

This is an advanced local-statistics result.

%%%%%%%%%%%%%%%%%%%%%%%%%%%%%%%%%%%%%%%%%%%%%%%%%%%%%%%%%%%%
\subsection{Universality}
\label{subsec:random-matrix-universality}
%%%%%%%%%%%%%%%%%%%%%%%%%%%%%%%%%%%%%%%%%%%%%%%%%%%%%%%%%%%%

Universality is the principle that many spectral statistics depend
primarily on broad structural features rather than the exact entry
distribution.

For example, matrices with quite different entry distributions may
share:

\[
\boxed{
\text{the same global spectral limit},
}
\]

\[
\boxed{
\text{the same local spacing law},
}
\]

or

\[
\boxed{
\text{the same edge fluctuation law}.
}
\]

This is one of the deepest themes in random matrix theory.

%%%%%%%%%%%%%%%%%%%%%%%%%%%%%%%%%%%%%%%%%%%%%%%%%%%%%%%%%%%%
\subsection{Sample Covariance Matrices}
\label{subsec:sample-covariance-random-matrix}
%%%%%%%%%%%%%%%%%%%%%%%%%%%%%%%%%%%%%%%%%%%%%%%%%%%%%%%%%%%%

Suppose

\[
X
\]

is a \(p\times n\) data matrix whose columns are observations.

A sample covariance-type matrix is

\[
\boxed{
S
=
\frac1n
XX^T.
}
\]

This is a \(p\times p\) positive semidefinite random matrix.

Its eigenvalues measure variability along principal directions.

%%%%%%%%%%%%%%%%%%%%%%%%%%%%%%%%%%%%%%%%%%%%%%%%%%%%%%%%%%%%
\subsection{Wishart Matrices}
\label{subsec:wishart-matrices}
%%%%%%%%%%%%%%%%%%%%%%%%%%%%%%%%%%%%%%%%%%%%%%%%%%%%%%%%%%%%

Suppose the columns of \(X\) are independent multivariate normal vectors

\[
\mathbf{x}_1,\ldots,\mathbf{x}_n
\sim
N_p(\mathbf{0},\Sigma).
\]

Then

\[
\boxed{
W
=
XX^T
}
\]

has a Wishart distribution.

One writes conceptually

\[
\boxed{
W
\sim
\operatorname{Wishart}_p(n,\Sigma).
}
\]

Wishart matrices are central in multivariate statistics.

%%%%%%%%%%%%%%%%%%%%%%%%%%%%%%%%%%%%%%%%%%%%%%%%%%%%%%%%%%%%
\subsection{Properties of Sample Covariance Matrices}
\label{subsec:sample-covariance-properties}
%%%%%%%%%%%%%%%%%%%%%%%%%%%%%%%%%%%%%%%%%%%%%%%%%%%%%%%%%%%%

Since

\[
S
=
\frac1nXX^T,
\]

we have

\[
\boxed{
S=S^T.
}
\]

Also,

\[
\boxed{
\mathbf{v}^TS\mathbf{v}
=
\frac1n
\|X^T\mathbf{v}\|^2
\geq0.
}
\]

Therefore \(S\) is positive semidefinite.

Its eigenvalues satisfy

\[
\boxed{
\lambda_i(S)\geq0.
}
\]

%%%%%%%%%%%%%%%%%%%%%%%%%%%%%%%%%%%%%%%%%%%%%%%%%%%%%%%%%%%%
\subsection{High-Dimensional Aspect Ratio}
\label{subsec:aspect-ratio-rmt}
%%%%%%%%%%%%%%%%%%%%%%%%%%%%%%%%%%%%%%%%%%%%%%%%%%%%%%%%%%%%

Classical statistics often assumes \(p\) fixed while

\[
n\to\infty.
\]

Random matrix theory studies the high-dimensional regime

\[
\boxed{
p,n\to\infty
}
\]

with

\[
\boxed{
\frac pn
\to
\gamma
}
\]

for some positive constant \(\gamma\).

The parameter \(\gamma\) is called the aspect ratio.

%%%%%%%%%%%%%%%%%%%%%%%%%%%%%%%%%%%%%%%%%%%%%%%%%%%%%%%%%%%%
\subsection{Marchenko--Pastur Law}
\label{subsec:marchenko-pastur-law}
%%%%%%%%%%%%%%%%%%%%%%%%%%%%%%%%%%%%%%%%%%%%%%%%%%%%%%%%%%%%

\begin{theorem}[Marchenko--Pastur Law]
For suitable random data matrices with standardized independent entries,
if

\[
\frac pn\to\gamma>0,
\]

then the empirical spectral distribution of

\[
\frac1nXX^T
\]

converges to the Marchenko--Pastur distribution.
\end{theorem}

For unit population variance, the continuous part has density

\[
\boxed{
\rho_{\mathrm{MP}}(x)
=
\frac{
\sqrt{
(b-x)(x-a)
}
}{
2\pi\gamma x
}
}
\]

for

\[
a\leq x\leq b,
\]

where

\[
\boxed{
a=(1-\sqrt{\gamma})^2,
}
\]

and

\[
\boxed{
b=(1+\sqrt{\gamma})^2.
}
\]

%%%%%%%%%%%%%%%%%%%%%%%%%%%%%%%%%%%%%%%%%%%%%%%%%%%%%%%%%%%%
\subsection{Marchenko--Pastur Atom at Zero}
\label{subsec:mp-atom-zero}
%%%%%%%%%%%%%%%%%%%%%%%%%%%%%%%%%%%%%%%%%%%%%%%%%%%%%%%%%%%%

When

\[
\gamma>1,
\]

the \(p\times p\) sample covariance matrix has rank at most \(n<p\).

Therefore it has at least

\[
p-n
\]

zero eigenvalues.

In the limiting Marchenko--Pastur law, this appears as a point mass at
zero of size

\[
\boxed{
1-\frac1\gamma.
}
\]

%%%%%%%%%%%%%%%%%%%%%%%%%%%%%%%%%%%%%%%%%%%%%%%%%%%%%%%%%%%%
\subsection{Interpretation of Marchenko--Pastur}
\label{subsec:interpret-marchenko-pastur}
%%%%%%%%%%%%%%%%%%%%%%%%%%%%%%%%%%%%%%%%%%%%%%%%%%%%%%%%%%%%

Even if the true covariance matrix is the identity,

\[
\Sigma=I,
\]

sample covariance eigenvalues do not all lie close to \(1\) when

\[
p
\]

is comparable to \(n\).

Instead, they spread across approximately

\[
\boxed{
\left[
(1-\sqrt{\gamma})^2,
(1+\sqrt{\gamma})^2
\right].
}
\]

This is a key high-dimensional phenomenon.

%%%%%%%%%%%%%%%%%%%%%%%%%%%%%%%%%%%%%%%%%%%%%%%%%%%%%%%%%%%%
\subsection{Worked Example: Marchenko--Pastur Support}
\label{subsec:worked-mp-support}
%%%%%%%%%%%%%%%%%%%%%%%%%%%%%%%%%%%%%%%%%%%%%%%%%%%%%%%%%%%%

\begin{workedexamplebox}{Aspect Ratio \(\gamma=1/4\)}

Suppose

\[
\gamma
=
\frac14.
\]

Then

\[
\sqrt{\gamma}
=
\frac12.
\]

Therefore,

\[
a
=
\left(
1-\frac12
\right)^2
=
\frac14,
\]

and

\[
b
=
\left(
1+\frac12
\right)^2
=
\frac94.
\]

\begin{answerbox}
\[
\boxed{
\text{Limiting spectral support}
=
\left[
\frac14,
\frac94
\right].
}
\]
\end{answerbox}

\end{workedexamplebox}

%%%%%%%%%%%%%%%%%%%%%%%%%%%%%%%%%%%%%%%%%%%%%%%%%%%%%%%%%%%%
\subsection{Singular Values}
\label{subsec:singular-values-random-matrices}
%%%%%%%%%%%%%%%%%%%%%%%%%%%%%%%%%%%%%%%%%%%%%%%%%%%%%%%%%%%%

For a rectangular matrix

\[
X\in\mathbb{R}^{p\times n},
\]

the singular values are

\[
s_1,\ldots,s_r,
\]

where

\[
r=\min(p,n).
\]

The squared singular values are the nonzero eigenvalues of

\[
XX^T
\]

and

\[
X^TX.
\]

Thus sample covariance spectra are directly connected to singular-value
spectra.

%%%%%%%%%%%%%%%%%%%%%%%%%%%%%%%%%%%%%%%%%%%%%%%%%%%%%%%%%%%%
\subsection{Singular Value Decomposition}
\label{subsec:random-svd}
%%%%%%%%%%%%%%%%%%%%%%%%%%%%%%%%%%%%%%%%%%%%%%%%%%%%%%%%%%%%

The singular value decomposition is

\[
\boxed{
X
=
U\Sigma V^T,
}
\]

where

\[
\Sigma
=
\operatorname{diag}
(
s_1,\ldots,s_r
)
\]

up to rectangular padding.

For random data matrices, the singular values are random variables.

Their asymptotic behavior determines numerical conditioning and
statistical signal strength.

%%%%%%%%%%%%%%%%%%%%%%%%%%%%%%%%%%%%%%%%%%%%%%%%%%%%%%%%%%%%
\subsection{Largest Singular Value}
\label{subsec:largest-singular-value}
%%%%%%%%%%%%%%%%%%%%%%%%%%%%%%%%%%%%%%%%%%%%%%%%%%%%%%%%%%%%

For standardized rectangular random matrices, the largest singular value
is approximately determined by the upper edge of the
Marchenko--Pastur spectrum.

Since eigenvalues of

\[
\frac1nXX^T
\]

lie approximately below

\[
(1+\sqrt{\gamma})^2,
\]

the largest singular value of

\[
\frac1{\sqrt{n}}X
\]

is approximately

\[
\boxed{
1+\sqrt{\gamma}.
}
\]

%%%%%%%%%%%%%%%%%%%%%%%%%%%%%%%%%%%%%%%%%%%%%%%%%%%%%%%%%%%%
\subsection{Smallest Singular Value}
\label{subsec:smallest-singular-value}
%%%%%%%%%%%%%%%%%%%%%%%%%%%%%%%%%%%%%%%%%%%%%%%%%%%%%%%%%%%%

For

\[
0<\gamma<1,
\]

the lower Marchenko--Pastur edge is

\[
(1-\sqrt{\gamma})^2.
\]

Thus the smallest singular value of

\[
\frac1{\sqrt{n}}X
\]

is approximately

\[
\boxed{
1-\sqrt{\gamma}.
}
\]

The smallest singular value is especially important for invertibility
and numerical conditioning.

%%%%%%%%%%%%%%%%%%%%%%%%%%%%%%%%%%%%%%%%%%%%%%%%%%%%%%%%%%%%
\subsection{Condition Number}
\label{subsec:random-condition-number}
%%%%%%%%%%%%%%%%%%%%%%%%%%%%%%%%%%%%%%%%%%%%%%%%%%%%%%%%%%%%

For a nonsingular square matrix,

\[
\boxed{
\kappa_2(A)
=
\frac{
s_{\max}(A)
}{
s_{\min}(A)
}.
}
\]

Random matrix theory studies how large this ratio is likely to be.

Condition numbers matter in numerical linear algebra because they
control sensitivity to perturbations.

%%%%%%%%%%%%%%%%%%%%%%%%%%%%%%%%%%%%%%%%%%%%%%%%%%%%%%%%%%%%
\subsection{Principal Component Analysis}
\label{subsec:rmt-pca}
%%%%%%%%%%%%%%%%%%%%%%%%%%%%%%%%%%%%%%%%%%%%%%%%%%%%%%%%%%%%

Principal component analysis uses eigenvectors of the sample covariance
matrix.

If

\[
S
=
\frac1nXX^T,
\]

then the eigenvectors corresponding to large eigenvalues identify
directions of high sample variance.

In high dimensions, random matrix theory helps distinguish

\[
\boxed{
\text{signal eigenvalues}
}
\]

from

\[
\boxed{
\text{noise eigenvalues}.
}
\]

%%%%%%%%%%%%%%%%%%%%%%%%%%%%%%%%%%%%%%%%%%%%%%%%%%%%%%%%%%%%
\subsection{Spiked Covariance Models}
\label{subsec:spiked-covariance}
%%%%%%%%%%%%%%%%%%%%%%%%%%%%%%%%%%%%%%%%%%%%%%%%%%%%%%%%%%%%

A simple population covariance model is

\[
\boxed{
\Sigma
=
I
+
\theta
\mathbf{u}\mathbf{u}^T,
}
\]

where

\[
\|\mathbf{u}\|=1
\]

and \(\theta>0\).

This adds one special high-variance direction to otherwise isotropic
noise.

Such a model is called a spiked covariance model.

%%%%%%%%%%%%%%%%%%%%%%%%%%%%%%%%%%%%%%%%%%%%%%%%%%%%%%%%%%%%
\subsection{Signal Versus Noise Eigenvalues}
\label{subsec:signal-noise-eigenvalues}
%%%%%%%%%%%%%%%%%%%%%%%%%%%%%%%%%%%%%%%%%%%%%%%%%%%%%%%%%%%%

In a spiked model, a sufficiently strong population spike may produce a
sample eigenvalue separated from the Marchenko--Pastur bulk.

A weak spike may remain hidden inside the noise spectrum.

Thus high-dimensional data may exhibit a detection threshold.

%%%%%%%%%%%%%%%%%%%%%%%%%%%%%%%%%%%%%%%%%%%%%%%%%%%%%%%%%%%%
\subsection{BBP Phase Transition Preview}
\label{subsec:bbp-transition}
%%%%%%%%%%%%%%%%%%%%%%%%%%%%%%%%%%%%%%%%%%%%%%%%%%%%%%%%%%%%

The Baik--Ben Arous--Péché phase transition describes a threshold at
which a population spike separates from the random covariance bulk.

Conceptually,

\[
\boxed{
\text{weak signal}
\longrightarrow
\text{absorbed into noise bulk},
}
\]

while

\[
\boxed{
\text{strong signal}
\longrightarrow
\text{outlier eigenvalue}.
}
\]

This phenomenon is important in modern high-dimensional statistics.

%%%%%%%%%%%%%%%%%%%%%%%%%%%%%%%%%%%%%%%%%%%%%%%%%%%%%%%%%%%%
\subsection{Random Graph Matrices}
\label{subsec:random-graph-matrices}
%%%%%%%%%%%%%%%%%%%%%%%%%%%%%%%%%%%%%%%%%%%%%%%%%%%%%%%%%%%%

A graph with \(n\) vertices has adjacency matrix

\[
A=(A_{ij}),
\]

where

\[
A_{ij}
=
\begin{cases}
1,&\text{vertices \(i\) and \(j\) are connected},\\
0,&\text{otherwise}.
\end{cases}
\]

For a random graph, the adjacency matrix is therefore a random matrix.

Spectral random graph theory studies the eigenvalues and eigenvectors of
this matrix.

%%%%%%%%%%%%%%%%%%%%%%%%%%%%%%%%%%%%%%%%%%%%%%%%%%%%%%%%%%%%
\subsection{Erd\H{o}s--R\'enyi Adjacency Matrix}
\label{subsec:erdos-renyi-adjacency}
%%%%%%%%%%%%%%%%%%%%%%%%%%%%%%%%%%%%%%%%%%%%%%%%%%%%%%%%%%%%

For the random graph

\[
G(n,p),
\]

each undirected edge appears independently with probability \(p\).

Therefore, for \(i<j\),

\[
\boxed{
A_{ij}
\sim
\operatorname{Bernoulli}(p).
}
\]

The matrix is symmetric:

\[
A_{ij}=A_{ji}.
\]

After centering,

\[
A-p(J-I),
\]

the matrix resembles a Wigner-type random matrix under suitable
regimes.

%%%%%%%%%%%%%%%%%%%%%%%%%%%%%%%%%%%%%%%%%%%%%%%%%%%%%%%%%%%%
\subsection{Graph Spectrum and Structure}
\label{subsec:graph-spectrum-random}
%%%%%%%%%%%%%%%%%%%%%%%%%%%%%%%%%%%%%%%%%%%%%%%%%%%%%%%%%%%%

Eigenvalues of graph matrices reveal information about:

\[
\boxed{
\text{connectivity},
}
\]

\[
\boxed{
\text{community structure},
}
\]

\[
\boxed{
\text{expansion},
}
\]

and

\[
\boxed{
\text{randomness}.
}
\]

Large eigenvalue outliers may indicate low-rank structure such as
communities.

%%%%%%%%%%%%%%%%%%%%%%%%%%%%%%%%%%%%%%%%%%%%%%%%%%%%%%%%%%%%
\subsection{Random Matrix Perturbations}
\label{subsec:random-matrix-perturbations}
%%%%%%%%%%%%%%%%%%%%%%%%%%%%%%%%%%%%%%%%%%%%%%%%%%%%%%%%%%%%

A common model is

\[
\boxed{
M
=
S+N,
}
\]

where

\[
S
\]

is a structured signal matrix and

\[
N
\]

is random noise.

The main questions are:

\[
\boxed{
\text{Can the signal eigenvalues be detected?}
}
\]

and

\[
\boxed{
\text{Can the signal eigenvectors be recovered?}
}
\]

This framework appears throughout statistics and machine learning.

%%%%%%%%%%%%%%%%%%%%%%%%%%%%%%%%%%%%%%%%%%%%%%%%%%%%%%%%%%%%
\subsection{Low-Rank Perturbations}
\label{subsec:low-rank-perturbations}
%%%%%%%%%%%%%%%%%%%%%%%%%%%%%%%%%%%%%%%%%%%%%%%%%%%%%%%%%%%%

If

\[
S
\]

has small rank, then

\[
M=S+N
\]

is a low-rank deformation of a random matrix.

Such models help explain:

\[
\text{PCA},
\]

\[
\text{matrix denoising},
\]

\[
\text{community detection},
\]

and

\[
\text{latent-factor models}.
\]

%%%%%%%%%%%%%%%%%%%%%%%%%%%%%%%%%%%%%%%%%%%%%%%%%%%%%%%%%%%%
\subsection{Resolvent}
\label{subsec:resolvent-random-matrix}
%%%%%%%%%%%%%%%%%%%%%%%%%%%%%%%%%%%%%%%%%%%%%%%%%%%%%%%%%%%%

For a matrix \(A\) and complex number \(z\) not in its spectrum, define
the resolvent

\[
\boxed{
R_A(z)
=
(A-zI)^{-1}.
}
\]

The resolvent is one of the most important analytic tools in random
matrix theory.

It packages spectral information into a matrix-valued analytic
function.

%%%%%%%%%%%%%%%%%%%%%%%%%%%%%%%%%%%%%%%%%%%%%%%%%%%%%%%%%%%%
\subsection{Normalized Trace of the Resolvent}
\label{subsec:normalized-resolvent-trace}
%%%%%%%%%%%%%%%%%%%%%%%%%%%%%%%%%%%%%%%%%%%%%%%%%%%%%%%%%%%%

The normalized trace is

\[
\boxed{
m_A(z)
=
\frac1n
\operatorname{tr}
(A-zI)^{-1}.
}
\]

Since the eigenvalues are \(\lambda_i\),

\[
\boxed{
m_A(z)
=
\frac1n
\sum_{i=1}^{n}
\frac1{\lambda_i-z}.
}
\]

This is the Stieltjes transform of the empirical spectral distribution,
up to the chosen sign convention.

%%%%%%%%%%%%%%%%%%%%%%%%%%%%%%%%%%%%%%%%%%%%%%%%%%%%%%%%%%%%
\subsection{Stieltjes Transform}
\label{subsec:stieltjes-transform-rmt}
%%%%%%%%%%%%%%%%%%%%%%%%%%%%%%%%%%%%%%%%%%%%%%%%%%%%%%%%%%%%

For a probability measure \(\mu\), define

\[
\boxed{
m_\mu(z)
=
\int_{\mathbb{R}}
\frac1{x-z}
\,d\mu(x),
}
\]

for

\[
z\in\mathbb{C}\setminus\mathbb{R}.
\]

If

\[
\mu=\mu_A,
\]

then

\[
m_\mu(z)
=
\frac1n
\operatorname{tr}(A-zI)^{-1}.
\]

Thus spectral convergence can be studied through analytic transforms.

%%%%%%%%%%%%%%%%%%%%%%%%%%%%%%%%%%%%%%%%%%%%%%%%%%%%%%%%%%%%
\subsection{Why Stieltjes Transforms Are Useful}
\label{subsec:why-stieltjes-transform}
%%%%%%%%%%%%%%%%%%%%%%%%%%%%%%%%%%%%%%%%%%%%%%%%%%%%%%%%%%%%

Stieltjes transforms convert a probability measure into a smooth
analytic function.

They are useful because:

\[
\boxed{
\text{matrix identities}
}
\]

can often be applied directly to resolvents.

After establishing convergence of the transforms, one can recover the
limiting spectral measure.

%%%%%%%%%%%%%%%%%%%%%%%%%%%%%%%%%%%%%%%%%%%%%%%%%%%%%%%%%%%%
\subsection{Semicircle Stieltjes Equation Preview}
\label{subsec:semicircle-stieltjes-equation}
%%%%%%%%%%%%%%%%%%%%%%%%%%%%%%%%%%%%%%%%%%%%%%%%%%%%%%%%%%%%

The Stieltjes transform of the standard semicircle law satisfies a
quadratic equation of the form

\[
\boxed{
m(z)^2+zm(z)+1=0,
}
\]

with the appropriate branch chosen so that

\[
m(z)
\sim
-\frac1z
\]

as

\[
|z|\to\infty.
\]

Solving gives

\[
\boxed{
m(z)
=
\frac{
-z+\sqrt{z^2-4}
}{2}
}
\]

with an appropriate complex branch convention.

%%%%%%%%%%%%%%%%%%%%%%%%%%%%%%%%%%%%%%%%%%%%%%%%%%%%%%%%%%%%
\subsection{Resolvent Identity}
\label{subsec:resolvent-identity}
%%%%%%%%%%%%%%%%%%%%%%%%%%%%%%%%%%%%%%%%%%%%%%%%%%%%%%%%%%%%

For matrices \(A\) and \(B\),

\[
\boxed{
(A-zI)^{-1}
-
(B-zI)^{-1}
=
(A-zI)^{-1}
(B-A)
(B-zI)^{-1}.
}
\]

This identity allows comparison of spectra under perturbations.

It is heavily used in modern random matrix proofs.

%%%%%%%%%%%%%%%%%%%%%%%%%%%%%%%%%%%%%%%%%%%%%%%%%%%%%%%%%%%%
\subsection{Rank-One Updates}
\label{subsec:rank-one-updates-rmt}
%%%%%%%%%%%%%%%%%%%%%%%%%%%%%%%%%%%%%%%%%%%%%%%%%%%%%%%%%%%%

Suppose

\[
A
\]

is perturbed by

\[
\mathbf{u}\mathbf{v}^T.
\]

Rank-one update formulas relate

\[
(A+\mathbf{u}\mathbf{v}^T)^{-1}
\]

to

\[
A^{-1}.
\]

The Sherman--Morrison identity gives

\[
\boxed{
(A+\mathbf{u}\mathbf{v}^T)^{-1}
=
A^{-1}
-
\frac{
A^{-1}\mathbf{u}\mathbf{v}^TA^{-1}
}{
1+\mathbf{v}^TA^{-1}\mathbf{u}
}
}
\]

when the denominator is nonzero.

Such identities are extremely useful for random covariance matrices.

%%%%%%%%%%%%%%%%%%%%%%%%%%%%%%%%%%%%%%%%%%%%%%%%%%%%%%%%%%%%
\subsection{Concentration of Spectral Quantities}
\label{subsec:spectral-concentration}
%%%%%%%%%%%%%%%%%%%%%%%%%%%%%%%%%%%%%%%%%%%%%%%%%%%%%%%%%%%%

Many random matrix statistics concentrate sharply around deterministic
values.

Examples include:

\[
\boxed{
\|A\|_2,
}
\]

\[
\boxed{
\frac1n\operatorname{tr}f(A),
}
\]

and

\[
\boxed{
\lambda_{\max}(A).
}
\]

Concentration inequalities allow random spectral behavior to be
controlled with high probability.

%%%%%%%%%%%%%%%%%%%%%%%%%%%%%%%%%%%%%%%%%%%%%%%%%%%%%%%%%%%%
\subsection{Matrix Concentration Preview}
\label{subsec:matrix-concentration-preview}
%%%%%%%%%%%%%%%%%%%%%%%%%%%%%%%%%%%%%%%%%%%%%%%%%%%%%%%%%%%%

Scalar concentration inequalities have matrix analogues.

Examples include:

\[
\boxed{
\text{matrix Bernstein inequality},
}
\]

\[
\boxed{
\text{matrix Chernoff inequality},
}
\]

and

\[
\boxed{
\text{matrix Hoeffding inequality}.
}
\]

These bound eigenvalues of sums of independent random matrices.

%%%%%%%%%%%%%%%%%%%%%%%%%%%%%%%%%%%%%%%%%%%%%%%%%%%%%%%%%%%%
\subsection{Random Matrix Sums}
\label{subsec:random-matrix-sums}
%%%%%%%%%%%%%%%%%%%%%%%%%%%%%%%%%%%%%%%%%%%%%%%%%%%%%%%%%%%%

Suppose

\[
Y_1,\ldots,Y_n
\]

are independent random matrices.

Then one may study

\[
\boxed{
S_n
=
\sum_{k=1}^{n}Y_k.
}
\]

Questions include:

\[
\boxed{
\|S_n-\mathbb{E}[S_n]\|_2
}
\]

and extreme eigenvalue deviations.

This is a noncommutative analogue of concentration for random sums.

%%%%%%%%%%%%%%%%%%%%%%%%%%%%%%%%%%%%%%%%%%%%%%%%%%%%%%%%%%%%
\subsection{Noncommutativity}
\label{subsec:rmt-noncommutativity}
%%%%%%%%%%%%%%%%%%%%%%%%%%%%%%%%%%%%%%%%%%%%%%%%%%%%%%%%%%%%

Unlike scalar random variables, matrices generally satisfy

\[
\boxed{
AB\neq BA.
}
\]

Therefore many classical scalar probability arguments cannot be applied
directly.

Random matrix probability combines probability with noncommutative
linear algebra.

This is one reason matrix concentration requires special methods.

%%%%%%%%%%%%%%%%%%%%%%%%%%%%%%%%%%%%%%%%%%%%%%%%%%%%%%%%%%%%
\subsection{Circular Law Preview}
\label{subsec:circular-law-preview}
%%%%%%%%%%%%%%%%%%%%%%%%%%%%%%%%%%%%%%%%%%%%%%%%%%%%%%%%%%%%

For non-Hermitian random matrices, eigenvalues may be complex.

A fundamental result is the circular law.

For suitable matrices with independent centered entries, the eigenvalues
of

\[
\frac1{\sqrt{n}}A_n
\]

become asymptotically distributed uniformly over the unit disk in the
complex plane.

Thus,

\[
\boxed{
\text{Hermitian matrices}
\longrightarrow
\text{semicircle-type real spectrum},
}
\]

while

\[
\boxed{
\text{non-Hermitian i.i.d. matrices}
\longrightarrow
\text{circular complex spectrum}.
}
\]

%%%%%%%%%%%%%%%%%%%%%%%%%%%%%%%%%%%%%%%%%%%%%%%%%%%%%%%%%%%%
\subsection{Ginibre Ensembles Preview}
\label{subsec:ginibre-ensembles}
%%%%%%%%%%%%%%%%%%%%%%%%%%%%%%%%%%%%%%%%%%%%%%%%%%%%%%%%%%%%

The Ginibre ensembles are non-Hermitian Gaussian random matrices.

For the complex Ginibre ensemble, entries are independent complex
Gaussian variables.

The resulting eigenvalues occupy a two-dimensional region in the
complex plane.

This contrasts strongly with GOE and GUE, where all eigenvalues are
real.

%%%%%%%%%%%%%%%%%%%%%%%%%%%%%%%%%%%%%%%%%%%%%%%%%%%%%%%%%%%%
\subsection{Free Probability Preview}
\label{subsec:free-probability-preview}
%%%%%%%%%%%%%%%%%%%%%%%%%%%%%%%%%%%%%%%%%%%%%%%%%%%%%%%%%%%%

Free probability is a noncommutative probability theory designed in
part to describe large random matrices.

Classical independence has a noncommutative analogue called freeness.

Under appropriate conditions, large independent random matrices can
behave asymptotically like free random variables.

%%%%%%%%%%%%%%%%%%%%%%%%%%%%%%%%%%%%%%%%%%%%%%%%%%%%%%%%%%%%
\subsection{Free Additive Convolution Preview}
\label{subsec:free-additive-convolution}
%%%%%%%%%%%%%%%%%%%%%%%%%%%%%%%%%%%%%%%%%%%%%%%%%%%%%%%%%%%%

In classical probability, the distribution of

\[
X+Y
\]

for independent variables is given by convolution.

In free probability, the spectral distribution of sums of asymptotically
free matrices is described by free additive convolution:

\[
\boxed{
\mu_{A+B}
\approx
\mu_A
\boxplus
\mu_B.
}
\]

The symbol

\[
\boxplus
\]

denotes free additive convolution.

%%%%%%%%%%%%%%%%%%%%%%%%%%%%%%%%%%%%%%%%%%%%%%%%%%%%%%%%%%%%
\subsection{Random Matrices and Number Theory}
\label{subsec:rmt-number-theory}
%%%%%%%%%%%%%%%%%%%%%%%%%%%%%%%%%%%%%%%%%%%%%%%%%%%%%%%%%%%%

Random matrix statistics also appear unexpectedly in number theory.

One famous connection compares local statistics of zeros of the Riemann
zeta function with eigenvalue statistics from random Hermitian matrices.

This does not mean the zeta zeros are literally random matrix
eigenvalues.

Rather, certain statistical patterns appear strikingly similar.

%%%%%%%%%%%%%%%%%%%%%%%%%%%%%%%%%%%%%%%%%%%%%%%%%%%%%%%%%%%%
\subsection{Random Matrices in Quantum Physics}
\label{subsec:rmt-quantum-physics}
%%%%%%%%%%%%%%%%%%%%%%%%%%%%%%%%%%%%%%%%%%%%%%%%%%%%%%%%%%%%

Random matrix theory originated partly in modeling complicated quantum
systems.

Instead of describing every microscopic interaction, one replaces the
Hamiltonian by a random matrix constrained by symmetry.

The resulting eigenvalues model energy levels.

Eigenvalue spacing statistics then describe spectral correlations.

%%%%%%%%%%%%%%%%%%%%%%%%%%%%%%%%%%%%%%%%%%%%%%%%%%%%%%%%%%%%
\subsection{Random Matrices in Wireless Communications}
\label{subsec:rmt-wireless}
%%%%%%%%%%%%%%%%%%%%%%%%%%%%%%%%%%%%%%%%%%%%%%%%%%%%%%%%%%%%

A multiple-input multiple-output communication channel may involve a
random matrix

\[
H.
\]

Channel capacity depends on quantities involving

\[
HH^*.
\]

Therefore eigenvalue distributions of random covariance-type matrices
determine large-system communication performance.

%%%%%%%%%%%%%%%%%%%%%%%%%%%%%%%%%%%%%%%%%%%%%%%%%%%%%%%%%%%%
\subsection{Random Matrices in Machine Learning}
\label{subsec:rmt-machine-learning}
%%%%%%%%%%%%%%%%%%%%%%%%%%%%%%%%%%%%%%%%%%%%%%%%%%%%%%%%%%%%

Random matrix ideas appear in:

\[
\boxed{
\text{high-dimensional covariance estimation},
}
\]

\[
\boxed{
\text{PCA},
}
\]

\[
\boxed{
\text{random features},
}
\]

\[
\boxed{
\text{neural-network Jacobians},
}
\]

and

\[
\boxed{
\text{spectral regularization}.
}
\]

Large-dimensional models often operate in regimes where classical
fixed-dimensional approximations fail.

%%%%%%%%%%%%%%%%%%%%%%%%%%%%%%%%%%%%%%%%%%%%%%%%%%%%%%%%%%%%
\subsection{Random Matrices and Numerical Linear Algebra}
\label{subsec:rmt-numerical-linear-algebra}
%%%%%%%%%%%%%%%%%%%%%%%%%%%%%%%%%%%%%%%%%%%%%%%%%%%%%%%%%%%%

Random matrices are used both as objects of analysis and as algorithmic
tools.

Applications include:

\[
\boxed{
\text{randomized low-rank approximation},
}
\]

\[
\boxed{
\text{random projections},
}
\]

\[
\boxed{
\text{preconditioning},
}
\]

and

\[
\boxed{
\text{sketching methods}.
}
\]

Spectral norm and singular-value estimates determine algorithmic
accuracy and stability.

%%%%%%%%%%%%%%%%%%%%%%%%%%%%%%%%%%%%%%%%%%%%%%%%%%%%%%%%%%%%
\subsection{Johnson--Lindenstrauss Connection Preview}
\label{subsec:jl-random-matrix-preview}
%%%%%%%%%%%%%%%%%%%%%%%%%%%%%%%%%%%%%%%%%%%%%%%%%%%%%%%%%%%%

Random projection matrices can approximately preserve pairwise
distances between many vectors.

A typical random linear map is

\[
\boxed{
x
\mapsto
\frac1{\sqrt{k}}Ax,
}
\]

where \(A\) has appropriately normalized random entries.

Concentration of random quadratic forms explains why distances are
approximately preserved.

This is closely connected to matrix concentration.

%%%%%%%%%%%%%%%%%%%%%%%%%%%%%%%%%%%%%%%%%%%%%%%%%%%%%%%%%%%%
\subsection{Random Matrix Modeling Workflow}
\label{subsec:rmt-modeling-workflow}
%%%%%%%%%%%%%%%%%%%%%%%%%%%%%%%%%%%%%%%%%%%%%%%%%%%%%%%%%%%%

When studying a random matrix problem:

\begin{enumerate}
    \item identify the matrix dimensions;
    \item determine whether the matrix is square or rectangular;
    \item identify symmetry or Hermitian structure;
    \item determine the entry distributions;
    \item identify independence assumptions;
    \item center the entries if needed;
    \item determine the correct normalization;
    \item decide whether eigenvalues or singular values are relevant;
    \item identify the empirical spectral distribution;
    \item determine whether a Wigner, covariance, graph, or
          non-Hermitian model applies;
    \item identify the large-dimensional regime;
    \item study the bulk spectrum;
    \item study extreme eigenvalues if needed;
    \item compare observed eigenvalues with the relevant random-noise
          baseline;
    \item interpret deviations as possible structure rather than
          automatically as signal.
\end{enumerate}

%%%%%%%%%%%%%%%%%%%%%%%%%%%%%%%%%%%%%%%%%%%%%%%%%%%%%%%%%%%%
\subsection{Common Errors}
\label{subsec:rmt-common-errors}
%%%%%%%%%%%%%%%%%%%%%%%%%%%%%%%%%%%%%%%%%%%%%%%%%%%%%%%%%%%%

\subsubsection{Ignoring Matrix Scaling}

For Wigner-type matrices, failing to divide by

\[
\sqrt{n}
\]

causes eigenvalues to grow with dimension.

The limiting semicircle law requires an appropriate normalization.

%%%%%%%%%%%%%%%%%%%%%%%%%%%%%%%%%%%%%%%%%%%%%%%%%%%%%%%%%%%%
\subsubsection{Assuming Entries and Eigenvalues Are Independent}

Even if matrix entries are independent, eigenvalues are generally
strongly dependent.

Eigenvalue repulsion is a major example.

%%%%%%%%%%%%%%%%%%%%%%%%%%%%%%%%%%%%%%%%%%%%%%%%%%%%%%%%%%%%
\subsubsection{Confusing Eigenvalues and Singular Values}

Eigenvalues may be negative or complex.

Singular values are always nonnegative.

For rectangular matrices, singular values are usually the more natural
spectral quantities.

%%%%%%%%%%%%%%%%%%%%%%%%%%%%%%%%%%%%%%%%%%%%%%%%%%%%%%%%%%%%
\subsubsection{Confusing \(XX^T\) with \(X^TX\)}

These matrices generally have different dimensions.

However, they share the same nonzero eigenvalues.

%%%%%%%%%%%%%%%%%%%%%%%%%%%%%%%%%%%%%%%%%%%%%%%%%%%%%%%%%%%%
\subsubsection{Forgetting the Aspect Ratio}

In high-dimensional covariance problems, spectral behavior depends on

\[
\gamma
=
\frac pn.
\]

The regime \(p\ll n\) behaves differently from \(p\) comparable to \(n\).

%%%%%%%%%%%%%%%%%%%%%%%%%%%%%%%%%%%%%%%%%%%%%%%%%%%%%%%%%%%%
\subsubsection{Assuming Sample Covariance Eigenvalues Equal Population Eigenvalues}

When dimension is large relative to sample size, sample eigenvalues are
systematically spread by random noise.

Marchenko--Pastur theory quantifies this effect.

%%%%%%%%%%%%%%%%%%%%%%%%%%%%%%%%%%%%%%%%%%%%%%%%%%%%%%%%%%%%
\subsubsection{Treating Every Large Eigenvalue as Signal}

Noise alone can generate large eigenvalues.

A suspected signal should be compared against an appropriate random
matrix null spectrum.

%%%%%%%%%%%%%%%%%%%%%%%%%%%%%%%%%%%%%%%%%%%%%%%%%%%%%%%%%%%%
\subsubsection{Assuming the Semicircle Law Applies to Every Random Matrix}

The semicircle law is associated with Wigner-type Hermitian models.

Sample covariance matrices instead lead to Marchenko--Pastur behavior.

Non-Hermitian i.i.d. matrices lead to circular-law behavior.

%%%%%%%%%%%%%%%%%%%%%%%%%%%%%%%%%%%%%%%%%%%%%%%%%%%%%%%%%%%%
\subsubsection{Ignoring Symmetry Constraints}

For a symmetric random matrix,

\[
A_{ij}=A_{ji}.
\]

Therefore the entries are not independent across the entire matrix.

Independence is imposed only on one triangular half under standard
models.

%%%%%%%%%%%%%%%%%%%%%%%%%%%%%%%%%%%%%%%%%%%%%%%%%%%%%%%%%%%%
\subsubsection{Assuming Global Laws Describe Extreme Eigenvalues Precisely}

The semicircle and Marchenko--Pastur laws describe global spectral
behavior.

Extreme eigenvalue fluctuations require separate edge theory.

%%%%%%%%%%%%%%%%%%%%%%%%%%%%%%%%%%%%%%%%%%%%%%%%%%%%%%%%%%%%
\subsubsection{Confusing Convergence of the ESD with Entrywise Convergence}

A spectral distribution may converge even though no individual matrix
entry converges.

These are fundamentally different types of limits.

%%%%%%%%%%%%%%%%%%%%%%%%%%%%%%%%%%%%%%%%%%%%%%%%%%%%%%%%%%%%
\subsection{Applications}
\label{subsec:rmt-applications}
%%%%%%%%%%%%%%%%%%%%%%%%%%%%%%%%%%%%%%%%%%%%%%%%%%%%%%%%%%%%

\begin{applicationbox}

\textbf{Multivariate Statistics.}
Wishart and Marchenko--Pastur theory describe high-dimensional covariance
matrices and principal components.

\medskip

\textbf{Machine Learning.}
Random matrix methods help analyze high-dimensional features, PCA,
spectral algorithms, and noisy low-rank models.

\medskip

\textbf{Network Science.}
Spectra of random adjacency and Laplacian matrices reveal connectivity,
communities, and departures from random-network structure.

\medskip

\textbf{Physics.}
Gaussian ensembles model complex quantum systems and energy-level
statistics.

\medskip

\textbf{Numerical Mathematics.}
Random matrix theory describes condition numbers, singular values, and
the behavior of randomized linear-algebra algorithms.

\medskip

\textbf{Wireless Communications.}
Large random channel matrices determine multi-antenna communication
capacity.

\medskip

\textbf{Finance.}
Random matrix filtering can be used to study noisy empirical covariance
matrices of many assets.

\medskip

\textbf{Number Theory.}
Random matrix eigenvalue statistics provide models for statistical
patterns observed among zeros of important analytic functions.

\end{applicationbox}

%%%%%%%%%%%%%%%%%%%%%%%%%%%%%%%%%%%%%%%%%%%%%%%%%%%%%%%%%%%%
\subsection{Exercises}
\label{subsec:rmt-exercises}
%%%%%%%%%%%%%%%%%%%%%%%%%%%%%%%%%%%%%%%%%%%%%%%%%%%%%%%%%%%%

\begin{exercise}[1]
Define a random matrix.
\end{exercise}

\begin{exercise}[1]
Define the empirical spectral distribution.
\end{exercise}

\begin{exercise}[1]
State the relationship between trace and eigenvalues.
\end{exercise}

\begin{exercise}[1]
State the relationship between determinant and eigenvalues.
\end{exercise}

\begin{exercise}[1]
Define a Wigner matrix.
\end{exercise}

\begin{exercise}[1]
State the Wigner semicircle law.
\end{exercise}

\begin{exercise}[1]
Define a sample covariance matrix.
\end{exercise}

\begin{exercise}[1]
Define the aspect ratio

\[
\gamma.
\]
\end{exercise}

\begin{exercise}[1]
State the Marchenko--Pastur support endpoints.
\end{exercise}

\begin{exercise}[1]
Define a Wishart matrix.
\end{exercise}

\begin{exercise}[1]
Define the resolvent of a matrix.
\end{exercise}

\begin{exercise}[1]
Define the Stieltjes transform of a probability measure.
\end{exercise}

\begin{exercise}[2]
Let

\[
A
=
\begin{pmatrix}
X&Y\\
Y&Z
\end{pmatrix}
\]

where \(X,Y,Z\) are random variables.

Explain why \(A\) is a random symmetric matrix.
\end{exercise}

\begin{exercise}[2]
Suppose \(A\) has eigenvalues

\[
-1,2,4.
\]

Compute

\[
\operatorname{tr}(A)
\]

and

\[
\det(A).
\]
\end{exercise}

\begin{exercise}[2]
For the same eigenvalues, compute

\[
\operatorname{tr}(A^2).
\]
\end{exercise}

\begin{exercise}[2]
Write the empirical spectral distribution for a matrix with eigenvalues

\[
1,1,3,5.
\]
\end{exercise}

\begin{exercise}[2]
For a standard Wigner model, explain why the normalization

\[
1/\sqrt{n}
\]

is natural.
\end{exercise}

\begin{exercise}[2]
Compute the first four Catalan numbers.
\end{exercise}

\begin{exercise}[2]
Verify that the odd moments of a symmetric semicircle density vanish.
\end{exercise}

\begin{exercise}[2]
Suppose

\[
\gamma=1.
\]

Find the Marchenko--Pastur support.
\end{exercise}

\begin{exercise}[2]
Suppose

\[
\gamma=\frac19.
\]

Find the Marchenko--Pastur support.
\end{exercise}

\begin{exercise}[2]
Suppose

\[
\gamma=4.
\]

Determine the limiting mass at zero in the Marchenko--Pastur law.
\end{exercise}

\begin{exercise}[2]
If

\[
X\in\mathbb{R}^{20\times100},
\]

what is the dimension of

\[
XX^T?
\]
\end{exercise}

\begin{exercise}[2]
For the same matrix, what is the dimension of

\[
X^TX?
\]
\end{exercise}

\begin{exercise}[2]
Explain why the nonzero eigenvalues of \(XX^T\) and \(X^TX\) agree.
\end{exercise}

\begin{exercise}[2]
Suppose

\[
\gamma=\frac14.
\]

Estimate the limiting largest singular value of

\[
X/\sqrt{n}
\]

under the standard covariance model.
\end{exercise}

\begin{exercise}[3]
Prove

\[
\operatorname{tr}(A^k)
=
\sum_i\lambda_i^k
\]

for a diagonalizable matrix.
\end{exercise}

\begin{exercise}[3]
Derive

\[
\frac1n\operatorname{tr}(A^k)
=
\int x^k\,d\mu_A(x).
\]
\end{exercise}

\begin{exercise}[3]
Expand

\[
\operatorname{tr}(A^4)
\]

as a sum over matrix entries.
\end{exercise}

\begin{exercise}[3]
Explain why centered independent entries cause many terms in expected
trace expansions to vanish.
\end{exercise}

\begin{exercise}[3]
Show that the standard semicircle density integrates to \(1\).
\end{exercise}

\begin{exercise}[3]
Compute the second moment of the standard semicircle distribution.
\end{exercise}

\begin{exercise}[3]
Compute the fourth moment and verify that it equals the Catalan number

\[
C_2=2.
\]
\end{exercise}

\begin{exercise}[3]
Show that a sample covariance matrix

\[
S=\frac1nXX^T
\]

is positive semidefinite.
\end{exercise}

\begin{exercise}[3]
Show that

\[
\operatorname{rank}(XX^T)
\leq
\min(p,n).
\]
\end{exercise}

\begin{exercise}[3]
Explain why \(\gamma>1\) forces zero eigenvalues in a \(p\times p\)
sample covariance matrix.
\end{exercise}

\begin{exercise}[3]
Show that detailed sample covariance behavior depends on the ratio

\[
p/n
\]

rather than only on \(p\) or \(n\) separately.
\end{exercise}

\begin{exercise}[3]
For a low-rank perturbation

\[
M=S+N,
\]

explain conceptually how a separated eigenvalue can reveal signal.
\end{exercise}

\begin{exercise}[3]
For a symmetric matrix, prove that the spectral norm equals

\[
\max_i|\lambda_i|.
\]
\end{exercise}

\begin{exercise}[3]
Derive the resolvent spectral representation

\[
\frac1n
\operatorname{tr}(A-zI)^{-1}
=
\frac1n
\sum_i
\frac1{\lambda_i-z}.
\]
\end{exercise}

\begin{exercise}[3]
Verify the resolvent identity

\[
(A-zI)^{-1}-(B-zI)^{-1}
=
(A-zI)^{-1}(B-A)(B-zI)^{-1}.
\]
\end{exercise}

\begin{exercise}[3]
Verify the Sherman--Morrison identity for an invertible matrix and a
rank-one perturbation under the required denominator condition.
\end{exercise}

\begin{exercise}[3]
Explain why eigenvalue repulsion is represented by a factor

\[
|\lambda_i-\lambda_j|^\beta.
\]
\end{exercise}

\begin{exercise}[3]
Compare Wigner, Wishart, and Ginibre ensembles.
\end{exercise}

\begin{exercise}[3]
Explain why the semicircle law and Marchenko--Pastur law describe
different random matrix structures.
\end{exercise}

\begin{exercise}[3]
Explain the difference between bulk spectral behavior and edge spectral
behavior.
\end{exercise}

\begin{exercise}[4]
Study a rigorous proof of the Wigner semicircle law using the moment
method.
\end{exercise}

\begin{exercise}[4]
Derive the Catalan-number combinatorics appearing in Wigner trace
moments.
\end{exercise}

\begin{exercise}[4]
Study a resolvent-based proof of the semicircle law.
\end{exercise}

\begin{exercise}[4]
Derive the Marchenko--Pastur law using Stieltjes transforms.
\end{exercise}

\begin{exercise}[4]
Study the joint eigenvalue density of GOE and GUE.
\end{exercise}

\begin{exercise}[4]
Investigate the role of the Vandermonde determinant in eigenvalue
repulsion.
\end{exercise}

\begin{exercise}[4]
Study sine-kernel universality in the bulk.
\end{exercise}

\begin{exercise}[4]
Investigate Tracy--Widom fluctuations of the largest eigenvalue.
\end{exercise}

\begin{exercise}[4]
Study the BBP phase transition in spiked covariance models.
\end{exercise}

\begin{exercise}[4]
Investigate the circular law for non-Hermitian random matrices.
\end{exercise}

\begin{exercise}[4]
Study the smallest singular value of random matrices.
\end{exercise}

\begin{exercise}[4]
Investigate matrix Bernstein inequalities.
\end{exercise}

\begin{exercise}[4]
Study spectra of Erd\H{o}s--R\'enyi random graphs.
\end{exercise}

\begin{exercise}[4]
Investigate community detection through spectral outliers in random
graph models.
\end{exercise}

\begin{exercise}[4]
Study free probability and asymptotic freeness.
\end{exercise}

\begin{exercise}[4]
Investigate free additive convolution.
\end{exercise}

\begin{exercise}[4]
Study random matrix models for high-dimensional PCA.
\end{exercise}

\begin{exercise}[4]
Investigate random matrices in wireless MIMO communication.
\end{exercise}

\begin{exercise}[4]
Study statistical similarities between random matrix eigenvalues and
zeros of the Riemann zeta function.
\end{exercise}

%%%%%%%%%%%%%%%%%%%%%%%%%%%%%%%%%%%%%%%%%%%%%%%%%%%%%%%%%%%%
\subsection{Conceptual Summary}
\label{subsec:rmt-summary}
%%%%%%%%%%%%%%%%%%%%%%%%%%%%%%%%%%%%%%%%%%%%%%%%%%%%%%%%%%%%

A random matrix is a matrix whose entries are random variables.

For a Hermitian matrix with eigenvalues

\[
\lambda_1,\ldots,\lambda_n,
\]

the empirical spectral distribution is

\[
\boxed{
\mu_A
=
\frac1n
\sum_{i=1}^{n}
\delta_{\lambda_i}.
}
\]

Spectral moments satisfy

\[
\boxed{
\int x^k\,d\mu_A(x)
=
\frac1n
\operatorname{tr}(A^k).
}
\]

For normalized Wigner matrices,

\[
\boxed{
\frac1{\sqrt{n}}W_n
}
\]

has an empirical spectral distribution approaching the semicircle law:

\[
\boxed{
\rho_{\mathrm{sc}}(x)
=
\frac1{2\pi}
\sqrt{4-x^2},
\qquad
|x|\leq2.
}
\]

The Gaussian ensembles exhibit eigenvalue repulsion through factors of
the form

\[
\boxed{
\prod_{i<j}
|\lambda_i-\lambda_j|^\beta.
}
\]

For sample covariance matrices,

\[
\boxed{
S
=
\frac1nXX^T,
}
\]

the high-dimensional parameter is

\[
\boxed{
\gamma
=
\frac pn.
}
\]

The Marchenko--Pastur support is

\[
\boxed{
\left[
(1-\sqrt{\gamma})^2,
(1+\sqrt{\gamma})^2
\right].
}
\]

The resolvent is

\[
\boxed{
R_A(z)
=
(A-zI)^{-1},
}
\]

and its normalized trace is the Stieltjes transform of the ESD:

\[
\boxed{
m_A(z)
=
\frac1n
\operatorname{tr}(A-zI)^{-1}.
}
\]

Random matrix theory therefore studies the transition

\[
\boxed{
\text{random entries}
\longrightarrow
\text{predictable large-scale spectral laws}
}
\]

while retaining subtle random fluctuations in individual eigenvalues.

Its central themes include

\[
\boxed{
\text{spectral limits}
+
\text{eigenvalue repulsion}
+
\text{universality}
+
\text{high-dimensional statistics}
+
\text{extreme eigenvalues}.
}
\]

%%%%%%%%%%%%%%%%%%%%%%%%%%%%%%%%%%%%%%%%%%%%%%%%%%%%%%%%%%%%
\subsection{Section Connections}
\label{subsec:rmt-connections}
%%%%%%%%%%%%%%%%%%%%%%%%%%%%%%%%%%%%%%%%%%%%%%%%%%%%%%%%%%%%

\chapterconnection{
Random Matrix Theory concludes the main progression of Chapter 7 by
extending probability from scalar random variables and stochastic
processes to random linear-algebraic structures. Joint distributions
describe random matrix entries, expectation and covariance describe
entry-level behavior, and limit theorems motivate the deterministic
spectral laws that emerge as matrix dimensions grow. The semicircle and
Marchenko--Pastur laws provide spectral analogues of classical
large-sample probability limits. Random matrices connect directly with
Linear Algebra, Combinatorics, Graph Theory, Functional Analysis,
Numerical Mathematics, Statistics, Machine Learning, Network Science,
Mathematical Physics, and Number Theory. Chapter 8 Statistics builds on
these probability foundations to study data collection, estimation,
inference, regression, experimental design, multivariate statistics,
Bayesian methods, time series, and statistical learning.
}

%%%%%%%%%%%%%%%%%%%%%%%%%%%%%%%%%%%%%%%%%%%%%%%%%%%%%%%%%%%%
% SECTION SOURCE TRACKING
%%%%%%%%%%%%%%%%%%%%%%%%%%%%%%%%%%%%%%%%%%%%%%%%%%%%%%%%%%%%

%------------------------------------------------------------
% 7.15 Random Matrix Theory
%------------------------------------------------------------
%
% Source basis:
%   - Standard random matrix theory.
%   - Standard high-dimensional covariance theory.
%   - Standard spectral probability methods.
%
% Used for:
%   - random matrix definitions;
%   - symmetric and Hermitian random matrices;
%   - empirical spectral distributions;
%   - traces and determinants;
%   - spectral norms;
%   - Wigner matrices;
%   - Wigner normalization;
%   - semicircle law;
%   - moment method;
%   - Catalan moments;
%   - Gaussian orthogonal ensemble;
%   - Gaussian unitary ensemble;
%   - Gaussian symplectic ensemble;
%   - beta ensembles;
%   - Vandermonde determinants;
%   - eigenvalue repulsion;
%   - bulk and edge spectra;
%   - largest eigenvalues;
%   - Tracy--Widom preview;
%   - local eigenvalue statistics;
%   - universality;
%   - sample covariance matrices;
%   - Wishart matrices;
%   - aspect ratios;
%   - Marchenko--Pastur law;
%   - zero eigenvalue mass;
%   - singular values;
%   - largest and smallest singular values;
%   - condition numbers;
%   - principal component analysis;
%   - spiked covariance models;
%   - BBP transition preview;
%   - random graph adjacency matrices;
%   - low-rank perturbations;
%   - resolvents;
%   - Stieltjes transforms;
%   - resolvent identities;
%   - rank-one updates;
%   - matrix concentration preview;
%   - circular law preview;
%   - Ginibre ensembles;
%   - free probability preview;
%   - free additive convolution;
%   - physics, number theory, communications, machine learning, and
%     numerical applications;
%   - exercises and chapter connections.
%
% Notes:
%   - This completes Chapter 7 Probability.
%   - High-dimensional statistical applications are revisited in
%     Chapter 8.
%   - Randomized numerical linear algebra is revisited in Chapter 11.
%   - Random graph spectral methods connect with Chapter 6.
%   - Free probability is included only as an advanced preview.
%
%%%%%%%%%%%%%%%%%%%%%%%%%%%%%%%%%%%%%%%%%%%%%%%%%%%%%%%%%%%%